\section{Where to Evolve?}
\label{sec:applications}

Self-evolving agents have facilitated advancements across a diverse array of domains and applications. 
Broadly, most of these applications can be systematically categorized into two groups: 
(1) general domain evolution, where agent systems evolve to expand their capabilities across a wide variety of tasks, mostly within the digital realm, and 
(2) specialized domain evolution, which evolves specifically to enhance their proficiency within particular task domains. 
In essence, evolution in general-purpose assistants focuses on transferring learned experience to a broader set of tasks, while evolution in specialized agents emphasizes deepening expertise within a specific domain.


\subsection{General Domain Evolution}

The first category, general domain evolution, refers to self-evolving agents designed for general-purpose applications, particularly as versatile digital assistants. These agents progressively enhance their capabilities to address a broad spectrum of user queries, especially in dynamic and diverse digital environments.
Technically speaking, these general assistant agent enhance their abilities primarily via three mechanisms: 
memory optimization, curriculum-driven training, and model-agent co‑evolution. 
These mechanisms collectively enable the agents to continuously adapt and effectively respond to increasingly complex user demands. 

\paragraph{Memory Mechanism.} The most common mechanism facilitating agent evolution is the memory mechanism, wherein agents summarize historical success/failure experiences~\citep{wang2023voyager, zhang2024offline} into memory representations~\citep{zhang2024survey}, anticipating that these distilled experiences will be beneficial when addressing previously unseen tasks.
For instance, Mobile-Agent-E~\citep{wang2025mobile} employs a long-term memory structure consisting of "Tips," which provide general guidelines, and "Shortcuts," representing reusable action sequences derived from past experiences. This self-evolutionary module supports the continuous enhancement of performance on complex smartphone tasks.
Another typical example is MobileSteward~\citep{liu2025mobilesteward}, which coordinates multiple app-specific Agents under a central Agent, with specialized modules for task scheduling, execution, and evaluation. It also incorporates a memory-based self-evolution mechanism that summarizes successful executions to improve future cross-app instruction handling.
Meanwhile, Generative Agents~\citep{park2023generative} store episodic memories of their experiences, synthesize higher-level reflections, and condition future planning on this self-reflection.
In these examples, memory serves as the foundation that enables agents to internalize past experiences, abstract high-level patterns, and refine their future behavior.

\begin{figure}[t]
  \centering
  \includegraphics[width=0.7\linewidth]{Figures/where.pdf}
  \caption{Categorization of where to evolve into two major types: General Domain Evolution, which focuses on broad capability enhancement across diverse tasks (e.g., memory mechanisms, co-evolution, curriculum training), and Specific Domain Evolution, which targets domain-specific expertise in areas such as coding, GUI, finance, medical, education, and others.}
  \label{fig:where}
\end{figure}


\paragraph{Model-Agent Co‑Evolution.} Another line of work is to perform Model-Agent Co‑evolution for LLM agents. 
UI‑Genie~\citep{xiao2025ui} constructs a specialized image-text reward model that scores trajectories at both step and task levels. It jointly fine-tunes the agent and reward model using synthetic trajectories—generated by controlled corruption and hard-negative mining—across multiple generations.
WebEvolver~\citep{fang2025webevolver} introduces a co-evolving world model LLM that simulates web environments. It generates synthetic training data by predicting next observations and enables look-ahead reasoning during inference, which greatly improves real-web task success. 
Absolute Zero~\citep{zhao2025absolute} co‑evolves a reasoning agent and its internal self‑reward model through reinforced self‑play. By adversarially generating increasingly challenging reasoning problems and optimizing the agent using internal self‑certainty as a reward signal, the framework simultaneously updates both the agent's policy and the self-rewarding mechanism.
Together, these methods demonstrate the effectiveness of co-evolving agents and auxiliary models (e.g., reward or world models) to achieve more robust, generalizable, and scalable learning in LLM agentic systems.

\paragraph{Curriculum-Driven Training.} 
Curriculum-driven training also serves as a critical mechanism for building a self-evolving general assistant. 
For example, WebRL~\citep{qi2024webrl} uses a self-evolving curriculum: when an agent fails, similar but manageable tasks are automatically generated. Coupled with a learned reward model and adaptive policy updates, this yields a success rate uplift on WebArena benchmarks. 
Voyager~\citep{wang2023voyager} similarly leverages an automatic, bottom‑up curriculum in Minecraft, where GPT‑4 proposes appropriate next tasks based on agent progress, building a growing code-based skill library through iterative prompting and environmental feedback. 
These approaches highlight how curriculum learning enables agents to autonomously expand their capabilities through iterative task adaptation.


\subsection{Specialized Domain Evolution}

In addition to general digital agents, self-evolving agents have also been effectively applied within specialized domains, where their evolution is tailored to significantly enhance performance within narrower task sets. 

\paragraph{Coding.}

The power of self-evolving agents extends directly to practical applications like coding, where their ability to autonomously adapt and improve offers a transformative approach to software development. SICA~\citep{robeyns2025self} demonstrates that a self-improving coding agent can autonomously edit its own codebase and improve its performance on benchmark tasks. EvoMAC~\citep{hu2024self} introduces a self-evolving paradigm on multi-agent collaboration networks, which automatically optimizes individual agent prompts and multi-agent workflows, significantly improving code generation performance by overcoming the limitations of manually designed systems. AgentCoder~\citep{huang2024agentcoder} also focuses on a multi-agent code generation framework that self-evolves through iterative refinement. A programmer agent continuously improves code based on feedback from a test executor agent, validated against independent test cases from a test designer, significantly boosting effectiveness and efficiency. Zhang et al.~\citep{zhang2025adaptive} enable LLM agents to continuously evolve by filtering high-quality answers, stratifying earned experiences by difficulty, and adaptively selecting demonstrations from self-generated data, leading to significant performance improvements and the construction of ML libraries. While these instances differ in their specific mechanisms—ranging from single-agent self-editing to complex multi-agent collaborative networks and experience-based learning—they commonly share the core principle of iterative self-improvement and autonomous adaptation to enhance coding capabilities. 
These advancements highlight how self-evolving agents can dramatically enhance coding efficiency and code quality by continuously learning and optimizing.

\paragraph{Graphical User Interfaces (GUI).}
Self‑evolving GUI agents extend LLM capabilities from pure text reasoning to direct manipulation of desktop, web, and mobile interfaces, where they must cope with large discrete action spaces, heterogeneous layouts, and partial visual observability. Yuan~\textit{et al.} couple pixel‑level vision with self‑reinforcement, enabling the agent to iteratively refine click–type grounding accuracy without additional human labels~\citep{yuan2025enhancing}.  On real desktop software, the Navi agent from \textit{WindowsAgentArena} replays and critiques its own failure trajectories, ultimately doubling its task‑completion rate across 150 Windows challenges~\citep{bonatti2024windows}.  For open‑web automation, \textit{WebVoyager} fuses screenshot features with chain‑of‑thought reflection; successive self‑fine‑tuning raises its end‑to‑end success on unseen sites from 30\,\% to 59\,\%~\citep{he2024webvoyager}, while ReAP adds episodic memories of past outcomes, recovering a further 29‑percentage‑point margin on previously failed queries~\citep{azam2025reap}.  Beyond RL and memory, \textit{AutoGUI} continuously mines functionality annotations from live interfaces to expand a reusable skill library each training cycle~\citep{li2025autogui}, and \textit{MobileUse} deploys a hierarchical self‑reflection stack that monitors, verifies, and revises smartphone actions in situ~\citep{wu2025mobileuse}.  Collectively, these systems epitomize the full triad of self‑evolution— what evolves (grounding modules, skill memories), when it evolves (offline consolidation vs.\ online reflection), and how it evolves (reinforcement learning, synthetic data, hierarchical monitoring)—charting a path toward universally competent interface agents. 



\paragraph{Financial.}

The primary bottleneck in customizing agents for specialized domains like financial tasks lies in efficiently constructing and integrating a domain-specific knowledge base into the agent’s learning process—a challenge that can be effectively mitigated by incorporating self-evolving mechanisms. 
QuantAgent~\citep{wang2024quantagent} proposed a two-layer framework that iteratively refines the agent's responses and automatically enhances its domain-specific knowledge base using feedback from simulated and real-world environments. This iterative process helps the agent progressively approximate optimal behavior, reduces reliance on costly human-curated datasets, and demonstrably improves its predictive accuracy and signal quality in trading tasks. 
TradingAgents~\citep{xiao2024tradingagents} incorporates dynamic processes such as reflection, reinforcement learning, and a feedback loop from real-world trading results, alongside collaborative debates, to continuously refine its strategies and enhance trading performance. 
These developments underscore the potential of self-evolving agents to revolutionize the financial domain by autonomously building domain expertise, adapting to dynamic market conditions, and continuously improving decision-making and trading performance.

\paragraph{Medical.}
Self-evolving agents have become a powerful paradigm in medical AI, where adaptability and the ability to evolve are essential for managing the complexity and ever-changing nature of real-world clinical practice.
One of the most prominent applications is hospital-scale simulation.
For example, Agent Hospital~\citep{agent2024hospital} creates closed environments with LLM-driven doctors, patients, and nurses, allowing the doctor agent to treat thousands of virtual cases. This process helps these agents autonomously refine and evolve their diagnostic strategies without manual labeling, ultimately achieving strong performance on USMLE-style exams.
Similarly, MedAgentSim~\citep{MedAgentSim2025} integrates an LLM doctor, patient, and tool agent. It records successful consultations as reusable trajectories and employs chain-of-thought reflection and consensus to drive self-evolution, improving success rates over successive interactions.
Another example is EvoPatient \citep{du2024evopatient} places a doctor agent and a patient agent in continuous dialogue. With each generation, they update their memory with high-quality exchanges: the patient develops more realistic symptom narratives, while the doctor learns to ask sharper questions. Notably, this happens without explicit gradient updates or hand-crafted rewards.
Reinforcement learning is also central to building adaptive medical agents. For instance, DoctorAgent-RL \citep{doctoragentRL2025} models consultations as a Markov decision process, using a reward function that scores diagnostic accuracy, coverage, and efficiency. This guides policy-gradient updates that help the agent ask more relevant questions and reach correct diagnoses faster than imitation-based approaches, thus achieving self-improvement.
In addition, automated architecture‑search approaches like~\textit{Learning to Be a Doctor} treat the workflow itself as an evolvable object, iteratively inserting specialist sub‑agents or new reasoning hops to cover observed failure modes and improve multimodal diagnostic accuracy~\citep{zhuang2025archsearch}. Finally, beyond clinical decision-making, self-evolving agents have also been extended to biomedical discovery. OriGene~\citep{zhang2025origene} functions as a virtual disease biologist that evolves by iteratively refining its analytical process. It leverages human and experimental feedback to update core reasoning templates, adjust tool usage strategies, and refine analytical protocols. Similarly, STELLA~\citep{jin2025stella} is a self-evolving biomedical research agent that improves over time by distilling successful reasoning workflows into reusable templates through its Template Library and expanding its Tool Ocean with external or newly assembled tools to meet emerging analytical needs.


\paragraph{Education.}

Self-evolving LLM agents have also found strong applications in the education domain.
At the learner level, self-evolving agents like the personalized tutor PACE~\citep{liu2025pace} adjust their prompts based on detailed student profiles and continually refine their questioning during conversations. Meanwhile, an LLM-to-LLM self-play framework generates diverse tutor–student dialogues that further fine-tune the agent, allowing its teaching strategies to evolve both during and after interactions.
Another example is MathVC~\citep{yue2025mathvc}, which employs symbolic persona profiles for virtual students and a meta-planner that orchestrates realistic problem-solving stages. This setup enables the agent’s conversational process to evolve step by step toward correct solutions, closely mirroring how collaborative learning naturally unfolds.
On the instructor side, self-evolving agent systems like the professional-development platform i‑vip~\citep{yang2025ivip} deploy a team of cooperating LLM agents—a coach, assessor, and feedback generator—that critique and enhance each other’s outputs in real time. These agents adapt their explanations based on teacher-learners’ responses and continue to evolve by incorporating expert feedback after deployment, thereby refining their prompt strategies over time
Similarly,  EduPlanner~\citep{zhang2025eduplanner} frames lesson-plan creation as an adversarial loop where a planner’s draft is repeatedly reviewed and refined by evaluator and optimizer agents until it meets diverse educational goals.
Similarly, \textsc{SEFL}~\citep{zhang2025sefl} uses teacher–student self-play to generate large sets of homework–feedback examples, which then fine-tune a lightweight feedback model. This self-evolving process significantly improves the clarity and usefulness of the comments.
Collectively, these examples illustrate how self-evolving LLM agents can dynamically adapt to both learners and instructors, driving more personalized, effective, and scalable educational experiences.


\paragraph{Others.}
Beyond the four major verticals discussed above, self-evolving agents demonstrate broader applicability, delivering superior adaptability and performance in specialized domains where conventional agents often fall short. 
For instance, Arxiv Copilot~\citep{lin2024arxivcopilot} learns and adapts by incorporating historical user interactions, including generated answers, research trends, and ideas, into its thought database, enhancing its ability to provide personalized and augmented academic assistance. 
In a very different context, Voyager~\citep{wang2023voyager}, an agent in the game Minecraft, excels at solving novel tasks from scratch in new worlds through a process of self-evolution. It continually refines its task goals via an automatic curriculum, expands its skill library, and enhances its actions using an iterative prompting mechanism without human intervention. Transitioning to domains that require explicit strategic planning, Agents-of-Change~\citep{belle2025agents} autonomously refines prompts and rewrites code based on iterative performance analysis and strategic research, thereby helping agents overcome inherent limitations in long-term strategic planning and achieve consistently superior and more coherent gameplay in complex environments like Settlers of Catan. Lastly, in the realm of diplomacy, Richelieu~\citep{zhao2024richelieu} introduces AI diplomacy agents that can self-evolve through their self-play mechanism, which allows the agent to augment its memory by acquiring diverse experiences without human data, thereby enhancing its strategic planning, reflection, and overall performance in diplomacy activities. While these diverse examples operate in distinct environments—from academic research and virtual game worlds to strategic board games and complex diplomatic negotiations—they all share the fundamental characteristic of leveraging continuous learning, self-refinement, and autonomous adaptation to achieve increasingly sophisticated and effective performance within their respective domains. 
These diverse examples reinforce the versatility of self-evolving agents, showcasing their growing potential to excel in a wide range of complex, dynamic, and human-like tasks beyond traditional domains.
